% TOP_PROBLEM: {"id": 3700010, "title": "Backward uniqueness for the heat equation", "category_id": 8, "category": {"id": 8, "name": "analysis", "display_name": "Analysis", "description": "Limits, continuity, calculus, and function theory.", "slug": "analysis", "order_index": 8, "created_at": "2026-07-31T15:26:25.671Z"}, "status": "partially_solved", "problem_number": "AMR-036-0010", "set_id": 13, "statusQualification": "Makes explicit that no spatial boundary condition is imposed. The source reports an unpublished equivalence claim, not an available proof."}
% FIRST_POSTED: 2026-10-01
% ENTRY_KIND: statement_only
% UnsolvedMath ID 3700010; source catalogue code AMR-036-0010
% !TeX program = lualatex
% Definition and sources only; this file is not an LLM solution attempt.
\documentclass[11pt,a4paper]{article}
\usepackage[margin=25mm]{geometry}
\usepackage{fontspec}
\setmainfont{lmroman10-regular.otf}[BoldFont=lmroman10-bold.otf,
 ItalicFont=lmroman10-italic.otf,BoldItalicFont=lmroman10-bolditalic.otf]
\usepackage{amsmath,amssymb,mathtools}
\usepackage{unicode-math}
\setmathfont{latinmodern-math.otf}
\AtBeginDocument{\let\setminus\smallsetminus}
\usepackage{xurl,hyperref}
\hypersetup{colorlinks=true,urlcolor=blue}
\setcounter{secnumdepth}{0}
\setlength{\parindent}{0pt}
\setlength{\parskip}{5pt}
\setlength{\emergencystretch}{3em}
\begin{document}
\section*{Backward uniqueness for the heat equation}\label{3700010}
{\small UnsolvedMath ID 3700010 · AMR-036-0010 · Analysis · AMR Open Problem Lists}
\subsection{Definitions and mathematical statement}
Let $D\subset\mathbb R^n$ be a connected open set, $n\geq1$,
whose boundary points are regular for the Dirichlet problem:
continuous boundary data are attained there by the corresponding
Perron harmonic solution. Compare the following two existence
properties of $D$.

Property I is the existence of a real-valued, nonzero, bounded
temperature $u$ on $D\times[0,1]$. More explicitly, $u$ is
continuous there, is twice differentiable in space and once in
time in the interior, and satisfies
\[
 u_t=\Delta_xu\quad\text{on }D\times(0,1),\qquad
 u(x,1)=0\quad(x\in D),\qquad
 \sup_{D\times[0,1]}|u|<\infty.
\]
Property II is the existence of a nonzero real harmonic
function $v$ on $D$ and constants $C,R>0$ such that
\[
             |v(x)|\leq C e^{-|x|^2}
       \quad(x\in D,\ |x|\geq R).
\]

Are Properties I and II equivalent for every such domain?
No value is prescribed for either function on the spatial
boundary $\partial D$; adding a homogeneous boundary
condition would change the question. Property II is a
uniform Gaussian decay condition along all directions to
infinity within $D$, and is vacuous at infinity when $D$
is bounded. Property I asks whether a bounded nontrivial
temperature can become identically zero at a later time.
The target is the full domain-dependent equivalence, beyond
special geometries such as intervals or cones.
\subsection{Short English statement}
Does failure of backward uniqueness for bounded heat solutions exactly correspond to a Gaussian-decaying harmonic function?
\subsection{Sources}
\begin{itemize}
\item[S1] ULAM.AI and the UnsolvedMath Contributors, supplied problems.json, record 3700010 (AMR-036-0010), AMR Open Problem Lists. Catalogue identity and source transcription; CC BY 4.0.
\url{https://huggingface.co/datasets/ulamai/UnsolvedMath}.
\item[S2] Eremenko - My favorite unsolved problems, Backward uniqueness problem. Original source indexed by AMR-036-0010.
\url{https://www.math.purdue.edu/~eremenko/uns1.html}.
\item[S3] A. Eremenko, Backward uniqueness for the heat equation. Author-maintained primary problem note.
\url{https://www.math.purdue.edu/~eremenko/dvi/heatproblem.pdf}.
\end{itemize}
\end{document}
